\begin{definition}[The algorithm of the round's graded-agreement composition]
  \label{def:aba-gbca-algorithm-by-abdy}
  \lean{PLTS.ABA.GBCA.ByABDY.boundOf}\lean{PLTS.ABA.GBCA.ByABDY.Algorithm}\lean{PLTS.ABA.GBCA.ByABDY.composition_projects}\leanok
  \uses{def:aba-labels, def:aba-parameters, def:aba-round-messages}
  The algorithm one graded-agreement round runs, as a relation on the state that round is composed over, with one transition per line of \algoref{alg:gbca-impl}{the GBCA implementation}, which is ABDY22's Algorithm~6 (D18).
  That state is the pair of the $n$ round records of Definition~\ref{def:aba-programs} and the round's network state of Definition~\ref{def:aba-gbca-instance-by-abdy}: the network is a component of the state rather than a field of it, and the four names $process$, $sent$, $received$ and $F$ used below are accessors on that pair.
  Each process runs a five-level exchange over the message levels of Definition~\ref{def:aba-round-messages}: it multicasts $\langle\mathrm{INPUT},b\rangle$ on a call and relays it on $f{+}1$ receipts, and each higher level is multicast into its own write-once field under an $n{-}f$ quorum at the level below and prior participation (D8), the levels from $\mathrm{BIND}$ down requiring in addition that the process's own send at the level below is already out.
  The $\mathrm{VOTE}$ level requires no own send: the $\mathrm{ECHO}$ it reads is multicast by an \textbf{upon} handler and may still be pending.
  A process returns $(b,2)$, $(b,1)$ or $(\bot,0)$ by the three wait-until cases of the last block, read as the if/elif/else chain that block ends in: a return carries the denials of the cases above it beside the receipts its own case names, and requires that the process has been called and its own $\langle\mathrm{ECHO5},v\rangle$ is out (D8).
  The network is set-based and authenticated (D5): each sender owns one idempotent message set, the adversary schedules delivery, and a corrupted sender may inject an arbitrary message.
  The state is the protocol's own data, so the specification's $excluded$ and $grade$ appear nowhere in it, and the one field beyond that data is the round's bound bit, held by the network state and read by no program (D29).
  $\mathrm{boundOf}$ computes that bit from the round's sent sets, the corrupted set and the outcome: a value-bearing outcome announces the value it hands out, and an outcome carrying no value announces the payload of a correct $\langle\mathrm{VOTE}, b\rangle$ sender, $1$ where there is none.
  A correct $\langle\mathrm{VOTE}, b\rangle$ sender holds an $n{-}f$ $\langle\mathrm{ECHO}, b\rangle$ receipt quorum and at most one bit carries such a quorum, so on a reachable state the two bit branches exclude one another.
  A return of outcome $out$ announces the bit already on record if the round has returned before and $\mathrm{boundOf}$'s bit otherwise, and writes it back, so a round announces one bit on all of its returns.
  The announced bit enters no guard of the algorithm and no field a program holds.
  The characterisation ties the algorithm to the composition of Definition~\ref{def:aba-gbca-instance-by-abdy}: at a label of the round's interface, every transition of the composition is the algorithm's transition at the specification label that interface label projects to, one step for one step.
  Full transition inventory: \leandecl{PLTS.ABA.GBCA.ByABDY.Algorithm}; the characterisation: \leandecl{PLTS.ABA.GBCA.ByABDY.composition_projects}.
\end{definition}
